\documentclass[11pt,a4paper]{article}
\usepackage[T1]{fontenc}
\usepackage[utf8]{inputenc}
\usepackage{lmodern,microtype}
\usepackage[margin=25mm]{geometry}
\usepackage{amsmath,amssymb,amsthm,mathtools}
\usepackage{booktabs,array,longtable}
\usepackage{xcolor,listings,fancyhdr}
\usepackage[hidelinks]{hyperref}
\hypersetup{
  pdftitle={A finite-difference measure converse and a Jacobi-symbol determinant counterexample},
  pdfauthor={Research manuscript prepared with ChatGPT for Keio},
  pdfsubject={Chasse--Csordas Conjecture 2.22 and Sun Conjecture 5.6(ii)}}
\pagestyle{fancy}
\fancyhf{}
\fancyhead[L]{\small Finite differences and Jacobi-symbol determinants}
\fancyhead[R]{\small 7 November 2026}
\fancyfoot[C]{\thepage}
\setlength{\headheight}{14pt}
\setlength{\emergencystretch}{2em}
\raggedbottom
\allowdisplaybreaks
\newtheorem{theorem}{Theorem}[section]
\newtheorem{lemma}[theorem]{Lemma}
\newtheorem{proposition}[theorem]{Proposition}
\newtheorem{corollary}[theorem]{Corollary}
\theoremstyle{remark}
\newtheorem{remark}[theorem]{Remark}
\newcommand{\F}{\mathbb F}
\newcommand{\Z}{\mathbb Z}
\newcommand{\jac}[2]{\left(\frac{#1}{#2}\right)}
\DeclareMathOperator{\Res}{Res}
\DeclareMathOperator{\diag}{diag}
\lstset{
  language=Python,basicstyle=\ttfamily\footnotesize,
  keywordstyle=\color{blue!45!black},
  commentstyle=\color{black!65},stringstyle=\color{black!75},
  showstringspaces=false,breaklines=true,
  columns=fullflexible,keepspaces=true,
  frame=single,rulecolor=\color{black!20},
  xleftmargin=4pt,xrightmargin=4pt,
  aboveskip=10pt,belowskip=10pt}

\DeclareMathOperator{\meas}{meas}
\DeclareMathOperator{\Gal}{Gal}

\title{\textbf{A finite-difference measure converse\\
and a Jacobi-symbol determinant\\counterexample}\\[7pt]
\large Chasse--Csordas Conjecture 2.22 and Sun Conjecture 5.6(ii)}
\author{Research manuscript prepared with ChatGPT for Keio}
\date{7 November 2026}
\begin{document}
\maketitle
\begin{abstract}
We give complete proofs concerning two previously stated conjectures.
First, if a real polynomial $p$ of degree $n\ge2$ satisfies
\[
 \meas\left\{x\in\mathbb R:
 \frac{p(x+1)-p(x)}{p(x)}\ge\lambda\right\}=\frac n\lambda
 \qquad(\lambda>0),
\]
then all its zeros are real and simple, and consecutive zeros are at
least one unit apart. This proves the implication in Chasse--Csordas
Conjecture~2.22. In fact, the identity on any nonempty open interval of
positive levels suffices. A rational measure-rigidity argument gives
positive simple real poles after cancellation; a divisor-count identity
then recovers the roots and multiplicities of the original polynomial.
The inverse-branch mechanism is an adaptation of Szeg\H{o}'s 1934 method,
and is not claimed as a historically new principle.

Second, for Sun's original truncated Jacobi-symbol determinant
$D_n=\{5,5\}_n$, we prove
\[
 D_{5469}\equiv2\pmod3,\qquad
 D_{5469}\equiv1089\pmod{1823},\qquad
 \jac{D_{5469}}{5469}=-1.
\]
Since $5469\equiv9\pmod{60}$ and $5469>69$, this is a counterexample to
Conjecture~5.6(ii). The proof uses a complementary inverse minor,
character diagonalization, and exact polynomial arithmetic. A complete
standard-library Python certificate is embedded in this single source.
Mathematical conclusions, historical attribution, and unverified
priority claims are kept separate.
\end{abstract}

\section{The finite-difference measure conjecture}
\label{sec:cc-statement}

For a real polynomial $p$ of degree $n\ge2$, define
\[
 F_p(x)=\frac{p(x+1)-p(x)}{p(x)}.
\]
Its domain excludes the real zeros of $p$. Removing or filling any
removable singularities changes a level set only at finitely many points,
so it does not change its Lebesgue measure.
For a polynomial with real zeros
$r_1\le\cdots\le r_n$, listed with multiplicities, write
\[
 \mu(p)=\min_{1\le j<n}(r_{j+1}-r_j).
\]
Thus $\mu(p)\ge1$ includes simplicity of every zero.

Chasse and Csordas \cite[Conjecture~2.22, p.~88]{ChasseCsordas2012}
ask whether the exact identity
\begin{equation}\label{eq:cc-original}
 \meas\{x\in\mathbb R:F_p(x)\ge\lambda\}=\frac n\lambda
 \quad\hbox{for every }\lambda>0
\end{equation}
forces $p$ to have only real zeros and $\mu(p)\ge1$.
No real-rootedness or simplicity hypothesis is imposed in this converse.
We prove it in the following slightly stronger form.

\begin{theorem}[Finite-difference measure characterization]
\label{thm:cc}
Let $p\in\mathbb R[x]$ have degree $n\ge2$. The following are equivalent.
\begin{enumerate}
\item All zeros of $p$ are real and $\mu(p)\ge1$.
\item Equation~\eqref{eq:cc-original} holds for every $\lambda>0$.
\item There is a nonempty open interval $I\subset(0,\infty)$ on which
\eqref{eq:cc-original} holds for every $\lambda\in I$.
\end{enumerate}
In particular, the implication in Chasse--Csordas Conjecture~2.22 holds
under its original hypotheses.
\end{theorem}

The forward implication is included for completeness and is not a new
claim: it is Theorem~88 of Chasse's 2011 dissertation
\cite[pp.~45--46]{ChasseThesis}; the converse appears there as
Conjecture~89. The main algebraic mechanism below has a substantial
precedent in Szeg\H{o}'s solution of P\'olya's Problem~107
\cite[pp.~19--20]{Szego1934}. We provide all details needed for the
present local level law and for the recovery of the original polynomial
after cancellation.

\section{A local rational measure-rigidity theorem}
\label{sec:rigidity}

\begin{theorem}[Local Boole-law rigidity]\label{thm:rigidity}
Let $F=A/B$ be a reduced nonconstant real rational function such that
\begin{equation}\label{eq:asymptotic}
 F(x)=\frac c x+O(x^{-2})\quad(x\longrightarrow\infty),\qquad c>0.
\end{equation}
If, on some nonempty open interval $I\subset(0,\infty)$,
\begin{equation}\label{eq:local-boole}
 \meas\{x\in\mathbb R:F(x)\ge\lambda\}=\frac c\lambda
 \qquad(\lambda\in I),
\end{equation}
then
\begin{equation}\label{eq:positive-poles}
 F(x)=\sum_{j=1}^{m}\frac{\rho_j}{x-b_j},\qquad
 b_1<\cdots<b_m\in\mathbb R,\quad
 \rho_j>0,\quad \sum_{j=1}^m\rho_j=c.
\end{equation}
In particular, every pole of the reduced rational function is simple
and real. Here $m=\deg B$.
\end{theorem}

\begin{proof}
Normalize $B$ to be monic. Equation~\eqref{eq:asymptotic} implies
$m=\deg B\ge1$, $\deg A=m-1$, and that the leading coefficient of
$A$ is $c$.

\paragraph{Step 1: an identity among inverse branches.}
The polynomial
\[
 L(X,T)=T B(X)-A(X)
\]
has only finitely many exceptional nonzero parameter values at which
its roots are not simple. Indeed, it is separable over $\mathbb C(T)$,
as also follows from the irreducibility argument in Step~2.
Choose a connected open subinterval $J\subset I$ avoiding these values.
Its $m$ roots can be labeled by analytic algebraic branches
$\alpha_1(T),\ldots,\alpha_m(T)$ on a complex neighborhood of a point
of $J$ and then continued along $J$. No root is a pole of $F$, since
$A$ and $B$ are coprime. Their real or nonreal status, and the crossing
direction of every real root, are fixed on $J$.

For $T>0$, the set $\{x:F(x)\ge T\}$ is bounded because $F(x)$ tends
to zero at both ends of the real axis. Its endpoints are real roots
of $F(x)=T$ and real poles of $F$. The poles contribute a constant
$C$. A downward crossing contributes its coordinate positively to
the length, an upward crossing negatively, and a nonreal root does
not contribute. Hence the assumed identity is
\begin{equation}\label{eq:signed-trace}
 C+\sum_{i=1}^m e_i\alpha_i(T)=\frac cT,\qquad
 e_i\in\{-1,0,1\},
\end{equation}
where the $e_i$ are constant on $J$.
This analytic identity is an identity of algebraic germs: its
difference vanishes on a real interval and therefore on the complex
neighborhood by the identity theorem. We may consequently apply
automorphisms of a splitting field of $L$ over $\mathbb C(T)$.

\paragraph{Step 2: transitivity excludes missing branches and wrong signs.}
The polynomial $T B(X)-A(X)$ is irreducible in $\mathbb C[X,T]$.
In any factorization, its degree one in $T$ forces one factor to
belong to $\mathbb C[X]$; that factor would divide both $A$ and $B$,
and is therefore constant. It is primitive as a polynomial in $X$
over $\mathbb C[T]$: a common nonconstant parameter factor would
give $A=t_0B$ for some $t_0$, contrary to their different degrees.
Gauss's lemma gives irreducibility over $\mathbb C(T)$, and
characteristic zero gives separability. Thus its splitting-field
Galois group $G$ acts transitively on the $m$ roots.

Let $s=\sum_i e_i$ and let $S_B$ be the sum of the roots of $B$,
with multiplicities. Average all $G$-conjugates of
\eqref{eq:signed-trace}. Each root occurs with average coefficient
$s/m$, so
\[
 C+\frac{s}{m}\sum_{i=1}^m\alpha_i(T)=\frac cT.
\]
By Vieta's formula,
\[
 \sum_{i=1}^m\alpha_i(T)=S_B+\frac cT.
\]
Comparison of the coefficients of $1/T$ gives $s=m$. Since every
$e_i\le1$, this forces
\[
 e_1=\cdots=e_m=1.
\]
Every level root on $J$ is therefore real and is a downward crossing.

\paragraph{Step 3: the poles are simple, real, and positive.}
Fix $T\in J$ and order the level roots
$\alpha_1<\cdots<\alpha_m$. At minus infinity $F-T<0$, whereas
immediately to the left of each $\alpha_i$ it is positive and
immediately to the right it is negative. Each of the $m$ disjoint
intervals
\[
 (-\infty,\alpha_1),\quad
 (\alpha_1,\alpha_2),\quad\ldots,\quad
 (\alpha_{m-1},\alpha_m)
\]
must contain a pole. Otherwise continuity would require an additional
level root, or an upward crossing, to change the sign from negative
to positive. Since $B$ has degree $m$, these exhaust all its zeros:
there is exactly one simple real pole in each interval and no others.
Across each such pole the sign changes from negative on the left to
positive on the right. The residue is therefore positive.
The rational function is proper, so its partial-fraction expansion is
\eqref{eq:positive-poles}; comparison at infinity gives
$\sum_j\rho_j=c$.
\end{proof}

\begin{remark}[Attribution of the mechanism]\label{rem:szego}
Szeg\H{o}'s 1934 converse for rational Lebesgue-measure preservers
already uses signed algebraic inverse branches, Vieta's formula,
irreducibility, and transitive monodromy to exclude a missing branch
or an incorrect sign \cite[pp.~19--20]{Szego1934}.
The Galois average above is an adaptation of that mechanism to a local
positive-tail identity with a prescribed coefficient at infinity.
We do not present the underlying principle as new, or infer worldwide
priority for the finite-difference application from a bounded search.
\end{remark}

\section{Recovering all polynomial zeros after cancellation}
\label{sec:recover}

\begin{proof}[Proof of Theorem~\ref{thm:cc}]
The implication (2)$\Rightarrow$(3) is immediate. We first prove
(3)$\Rightarrow$(1).

For any real polynomial of degree $n$,
\[
 F_p(x)=\frac{p(x+1)-p(x)}{p(x)}
       =\frac nx+O(x^{-2}).
\]
Reduce this rational function before applying
Theorem~\ref{thm:rigidity}. The denominator degree $m$ after reduction
may be smaller than $n$, but the coefficient at infinity remains $n$.
We obtain
\[
 R(x):=\frac{p(x+1)}{p(x)}=1+F_p(x)
 =1+\sum_{j=1}^m\frac{\rho_j}{x-b_j},\qquad \rho_j>0.
\]

\paragraph{No nonreal root can be lost through all cancellations.}
Suppose that $p$ has a nonreal root. Among the roots on one of the
corresponding horizontal lines, choose $z$ with maximal real part.
Then $p(z+1)\ne0$, so $R$ has a nonreal pole at $z$.
This contradicts the real-pole conclusion of
Theorem~\ref{thm:rigidity}. Hence all roots of $p$ are real.

\paragraph{The reduced quotient has strictly alternating zeros and poles.}
For real $x$ away from the poles,
\[
 R'(x)=-\sum_{j=1}^m\frac{\rho_j}{(x-b_j)^2}<0.
\]
Its limit at both infinities is $1$, and its limits at each pole
are $-\infty$ from the left and $+\infty$ from the right.
It follows that $R$ has exactly one simple zero to the left of $b_1$
and one in each interval $(b_j,b_{j+1})$, with none to the right of
$b_m$. These $m$ zeros exhaust its numerator degree. Denoting them
by $a_1,\ldots,a_m$, their order is
\begin{equation}\label{eq:alternation}
 a_1<b_1<a_2<b_2<\cdots<a_m<b_m.
\end{equation}

\paragraph{A divisor count recovers multiplicity and spacing.}
Let $N_p(t)$ count the zeros of $p$ in $(-\infty,t]$, with their
original multiplicities. The zeros of $p(x+1)$ are the zeros of $p$
shifted left by one. Common factors cancel equally in the numerator
and denominator, so for every real $t$,
\begin{equation}\label{eq:divisor-count}
 N_p(t+1)-N_p(t)
 =\#\{j:a_j\le t\}-\#\{j:b_j\le t\}\in\{0,1\}.
\end{equation}
The final inclusion follows from \eqref{eq:alternation}.
If $p$ had a multiple root, or two distinct roots less than one apart,
we could place at least two roots, counted with multiplicities,
strictly inside some interval $(t,t+1)$. The left side of
\eqref{eq:divisor-count} would then be at least two, a contradiction.
All roots are consequently simple and $\mu(p)\ge1$.

\paragraph{The known forward implication.}
Suppose now that $r_1<\cdots<r_n$ are the roots of $p$ and
$r_{j+1}-r_j\ge1$. Partial fractions give
\[
 F_p(x)=\sum_{j=1}^n\frac{\rho_j}{x-r_j},\qquad
 \rho_j=\frac{p(r_j+1)}{p'(r_j)}
       =\prod_{k\ne j}\frac{r_j+1-r_k}{r_j-r_k}\ge0.
\]
Every factor is nonnegative by the spacing assumption. Zero residues
represent removable poles; discard them, retaining poles
$b_1<\cdots<b_m$ with positive residues. Their sum is $n$ by the
coefficient at infinity.
For each $\lambda>0$, strict decrease between poles shows that the
superlevel set is, up to finitely many endpoints,
\[
 \bigcup_{j=1}^m(b_j,\alpha_j],
\]
where $\alpha_j$ is the unique level root to the right of $b_j$
and before the next pole, or on $(b_m,\infty)$ for $j=m$.
Vieta's formula for $\lambda B-A$ gives
\[
 \meas\{F_p\ge\lambda\}
 =\sum_j(\alpha_j-b_j)=\frac n\lambda.
\]
This proves (1)$\Rightarrow$(2), and finishes the proof.
\end{proof}

\begin{remark}[Why reduction cannot be skipped]
For $p(x)=\prod_{j=0}^{n-1}(x-a-j)$,
\[
 F_p(x)=\frac{n}{x-a-(n-1)}.
\]
The reduced denominator has degree one, although $\deg p=n$.
This example attains $\mu(p)=1$ and obeys the exact measure law.
Conversely, for $p(x)=x^2(x-1)$ one obtains
\[
 F_p(x)=-\frac1x+\frac4{x-1},\qquad
 \meas\{F_p\ge1\}=2\sqrt5-1\ne3.
\]
The negative residue records the failure of the required root spacing.
These examples illustrate why the proof must retain the original
divisor multiplicities after reducing the quotient.
\end{remark}

\section{The unchanged Sun conjecture and its counterexample}\label{sec:sun-target}

Sun defines, for an odd integer \(n>1\),
\begin{equation}\label{eq:original}
\{c,d\}_n=
\det\!\left[\jac{j^2+cjk+dk^2}{n}\right]_{1<j,k<n-1},
\end{equation}
where the entries are Jacobi symbols. The indices are all integers
\(2,\ldots,n-2\), without a coprimality restriction.
Conjecture~5.6(ii) in \cite[p.~23]{Sun2024} asserts that
\(\jac{\{5,5\}_n}{n}=0\) whenever \(n\equiv11,19\pmod{20}\), or
whenever \(n\equiv9\pmod{60}\) and \(n>69\).
The assumption that \(n\) is a sum of two squares belongs to part~(i)
and is absent from part~(ii).

Li and Yan \cite[Theorem~1.1(f), Remark~1.2]{LiYan2026} prove the
\(11,19\pmod{20}\) cases of part~(ii). Their stated results leave the
\(9\pmod{60}\) clause untreated. The counterexample below addresses
exactly that clause.

\begin{theorem}\label{thm:main}
Set \(p=1823\), \(n=3p=5469\), and \(D_n=\{5,5\}_n\), with the
definition \eqref{eq:original}. Then
\begin{equation}\label{eq:main-residues}
D_n\equiv2\pmod3,\qquad D_n\equiv1089\pmod p.
\end{equation}
Equivalently, \(D_n\equiv2912\pmod{5469}\). Consequently \(\gcd(D_n,n)=1\) and
\[
\jac{D_n}{n}=-1.
\]
In particular, Sun's Conjecture~5.6(ii), as stated in
\cite{Sun2024}, is false.
\end{theorem}

The modulus \(n=5469\) is odd, exceeds \(69\), and is congruent to
\(9\) modulo \(60\). Since \(1089=33^2\) is nonzero modulo \(1823\),
\eqref{eq:main-residues} gives
\[
\jac{D_n}{3p}
=\jac{2}{3}\jac{33^2}{1823}
=(-1)(1)=-1.
\]
Thus proving the two displayed congruences suffices. No exact decimal
expansion of the large integer \(D_n\) is needed.

\paragraph{What is being proved computationally.}
This is a finite, exact certificate for a counterexample to a universal
assertion. It is not an extrapolation from a finite range of successful
tests. The matrices, fields, loop bounds, polynomial identities and
final residues are all specified below. The certificate uses integer
arithmetic throughout.

\section{Restoring the deleted indices}

Write \(q(X)=X^2+5X+5\). For an odd prime \(r\), define the full matrix
\[
A_r=\left[\jac{j^2+5jk+5k^2}{r}\right]_{j,k\in\F_r}.
\]
The row and column orders are the same. Define \(A_n\) by the same full matrix with Jacobi entries when \(n\) is composite. If \(n=3p\), with \(p\ne3\),
the Chinese remainder bijection
\(\Z/n\Z\longrightarrow\F_3\times\F_p\) gives
\begin{equation}\label{eq:tensor}
A_n\sim A_3\otimes A_p.
\end{equation}
Here \(\sim\) is simultaneous permutation of rows and columns, so it
preserves determinants. Equality of entries follows from
\(\jac{a}{3p}=\jac{a}{3}\jac{a}{p}\).
The determinant \(D_n\) is obtained from \(A_n\) by deleting the
three rows and columns indexed by
\[
\mathcal E=\{0,1,-1\}.
\]
Under the Chinese remainder ordering these indices become
\((0,0),(1,1),(-1,-1)\), respectively.

\begin{lemma}[Complementary principal minor]\label{lem:complement}
If \(M\) is an invertible matrix over a field and \(I\) is a set of
indices, then
\[
\det M_{I^c,I^c}=\det M\;\det(M^{-1})_{I,I}.
\]
\end{lemma}
\begin{proof}
After a simultaneous permutation, partition \(M\) into blocks indexed
by \(I^c,I\). If its upper-left block is invertible, block elimination
and the Schur complement prove the formula. Multiplying by suitable
powers of \(\det M\) makes the identity polynomial in the matrix
entries. The generic identity therefore holds for every invertible
\(M\), including cases where the upper-left block is singular.
\end{proof}

A direct calculation gives
\begin{equation}\label{eq:A3}
A_3=
\begin{pmatrix}
0&-1&-1\\
1&-1&1\\
1&1&-1
\end{pmatrix},
\qquad
\det A_3=-4,\qquad
A_3^{-1}=
\begin{pmatrix}
0&1/2&1/2\\
-1/2&-1/4&1/4\\
-1/2&1/4&-1/4
\end{pmatrix}.
\end{equation}
These formulas are valid over each field used below.

\begin{proposition}[Three-coordinate reduction]\label{prop:minor}
Let \(K\) be a field of odd characteristic, let \(p\ne3\) be an odd prime, and suppose
\(A_p\) is invertible over \(K\). If the entries of \(B=A_p^{-1}\)
at \(0,1,-1\) have the form
\begin{equation}\label{eq:Bsmall}
B_{\mathcal E,\mathcal E}
=
\begin{pmatrix}
0&u&u\\
-u&a&b\\
-u&b&a
\end{pmatrix},
\end{equation}
then, as an identity in \(K\),
\begin{equation}\label{eq:master}
D_{3p}=(-4)^p(\det A_p)^3\,
\frac{u^2(a+b)}8.
\end{equation}
\end{proposition}
\begin{proof}
By \eqref{eq:tensor},
\(\det A_{3p}=(-4)^p(\det A_p)^3\) and
\(A_{3p}^{-1}=A_3^{-1}\otimes B\) in the corresponding order.
Its principal submatrix at the three deleted indices is
\[
C=
\begin{pmatrix}
0&u/2&u/2\\
u/2&-a/4&b/4\\
u/2&b/4&-a/4
\end{pmatrix}.
\]
Expansion gives \(\det C=u^2(a+b)/8\).
Lemma~\ref{lem:complement} proves \eqref{eq:master}.
\end{proof}

Formula \eqref{eq:master} shows the necessary additional condition:
invertibility of the full local matrices alone does not ensure
nonvanishing after truncation. The inverse entries must also satisfy
\(a+b\ne0\). This is the cancellation that the certificate will check.
Restoration, character diagonalization and Chinese remainder
decomposition are established tools; see in particular
\cite[\S1.2, Lemmas~2.1 and~2.4]{LiYan2026}.
The complementary inverse minor is used here to prove nonvanishing
for the original truncated determinant.

\section{The congruence modulo 1823}\label{sec:sun-prime}

We first work over \(\F_p\), for a prime \(p>3\) with \(p\equiv3\pmod{20}\).
Put \(h=(p-1)/2\), and define
\begin{equation}\label{eq:coeffs}
q(X)^h=\sum_{s=0}^{p-1}a_sX^s\quad\text{in }\F_p[X].
\end{equation}
Quadratic reciprocity gives \(\jac5p=-1\), so \(a_0=-1\) and
\(a_{p-1}=1\).

\begin{lemma}\label{lem:prime}
Suppose \(a_s\ne0\) for \(1\le s\le p-2\). Define
\[
P=\prod_{s=1}^{p-2}a_s,\qquad
E=\sum_{\substack{1\le s\le p-2\\s\ {\rm even}}}a_s^{-1}.
\]
Then
\begin{equation}\label{eq:prime-formula}
\det A_p=P,\qquad D_{3p}=-P^3E\quad\text{in }\F_p.
\end{equation}
\end{lemma}
\begin{proof}
Let \(e_0\) be the vector supported at zero, and let \(v_0\) be
zero at zero and one at every nonzero element of \(\F_p\).
For \(1\le s\le p-2\), let \(v_s(0)=0\) and \(v_s(j)=j^s\)
for \(j\ne0\).
The vectors \(e_0,v_0,v_1,\ldots,v_{p-2}\) form a basis.
For nonzero \(j,k\), Euler's criterion gives
\[
(A_p)_{j,k}=q(j/k)^h=\sum_{t=0}^{p-1}a_tj^tk^{-t}.
\]
Using \(\sum_{k\ne0}k^r=-1\) when \(p-1\mid r\) and zero
otherwise, we obtain
\[
A_pv_s=-a_sv_s\quad(1\le s\le p-2),\qquad
A_pe_0=v_0,\qquad A_pv_0=e_0.
\]
The last equality uses \(a_0+a_{p-1}=0\) and
\(-(p-1)=1\) in \(\F_p\).
The two-dimensional block on \(e_0,v_0\) has determinant \(-1\).
There are \(p-2\) remaining, hence an odd number of, eigenvalues.
It follows that
\[
\det A_p=-\prod_{s=1}^{p-2}(-a_s)=P.
\]

For a nonzero index \(k\), the coordinate vector satisfies
\[
e_k=-v_0-\sum_{s=1}^{p-2}k^{-s}v_s.
\]
Applying \(A_p^{-1}\) gives
\[
A_p^{-1}e_k=-e_0+\sum_{s=1}^{p-2}a_s^{-1}k^{-s}v_s.
\]
Also \(A_p^{-1}e_0=v_0\). Thus \eqref{eq:Bsmall} holds with
\[
u=-1,\quad
a=\sum_{s=1}^{p-2}a_s^{-1},\quad
b=\sum_{s=1}^{p-2}(-1)^sa_s^{-1},
\qquad a+b=2E.
\]
Proposition~\ref{prop:minor} and \((-4)^p=-4\) in \(\F_p\)
now give \(D_{3p}=(-4)P^3(2E)/8=-P^3E\).
\end{proof}

\subsection{Exact specialization}
The complete calculation in Appendix~\ref{app:code} gives, for \(p=1823\),
\begin{equation}\label{eq:prime-data}
\begin{gathered}
a_s\ne0\quad(1\le s\le1821),\\
P=697,\qquad E=1662,\qquad
a=921,\quad b=580.
\end{gathered}
\end{equation}
In particular, \(a+b=1501=2E\) in \(\F_{1823}\), and the
complementary inverse determinant is
\[
\frac E4=1327.
\]
The full determinant is
\((-4)^{1823}697^3=641\) in this field. Therefore
\[
D_{5469}=641\cdot1327
=-697^3\cdot1662
=1089\quad\text{in }\F_{1823}.
\]
The program constructs every coefficient in \eqref{eq:coeffs}
by \(911\) multiplications by \(X^2+5X+5\).
It checks every required nonzero coefficient before taking any inverse.
The value \(E=1662\ne0\), in addition to \(P\ne0\), is the
operative noncancellation certificate.

\section{The congruence modulo 3}\label{sec:sun-three}

The prime local matrix \(A_p\) can also be studied over a field whose
characteristic differs from \(p\). We give all the details needed
for \(p=1823\) and \(K=\F_3\).
Set \(m=p-1=1822\). The number \(5\) generates \(\F_p^\times\).
Indeed, \(m=2\cdot911\), both \(p\) and \(911\) are prime, and
the program verifies \(5^{911}\not\equiv1\pmod p\) and
\(5^2\not\equiv1\pmod p\).

Order the nonzero indices as \(1,5,\ldots,5^{m-1}\), and define
\begin{equation}\label{eq:circulant}
f_k=\jac{q(5^k)}p,\qquad
F(X)=\sum_{k=0}^{m-1}f_kX^k,\qquad
S(X)=1+X+\cdots+X^{m-1},
\end{equation}
with the coefficients of \(F\) reduced modulo \(3\).
The nonzero-index block of \(A_p\) is the circulant
\((f_{i-j})_{0\le i,j<m}\), where subscripts are taken modulo \(m\).

Its row sum is zero already over the integers. To see this,
the discriminant of \(q\) is \(5\ne0\pmod p\), and
\(\sum_{x\in\F_p}\jac{q(x)}p=-1\).
The latter identity follows by completing the square and counting
solutions of \(y^2=x^2-d\), \(d\ne0\): the bijection
\((x-y,x+y)\) identifies these solutions with pairs of nonzero
elements with product \(d\), of which there are \(p-1\).
Removing \(x=0\), whose symbol is \(-1\), gives a zero row sum.
In particular \(F(1)=0\).

Since \(3\nmid m\), we have the direct sum of the constant vectors
and the subspace of nonzero-index vectors with coordinate sum zero.
On the two-dimensional space generated by \(e_0,v_0\), the full
matrix has the block
\[
\begin{pmatrix}0&-m\\1&0\end{pmatrix},
\]
whose determinant is \(m\).
On the complementary space, the circulant has eigenvalues
\(F(\zeta)\), where \(\zeta\) runs over the nontrivial \(m\)-th
roots of unity in an algebraic closure of \(\F_3\).
The inversion of \(\zeta\) merely permutes the eigenvalues.
Consequently
\begin{equation}\label{eq:resultant}
\det A_p=m\,\Res(S,F).
\end{equation}

Assume \(F\) is invertible modulo \(S\).
Compute \(H_0\), of degree at most \(m-2\), such that
\(FH_0\equiv1\pmod S\), and set
\begin{equation}\label{eq:H}
H=H_0-\frac{H_0(1)}mS.
\end{equation}
Then \(H(1)=0\), and
\[
FH\equiv1-\frac1mS\pmod{X^m-1}.
\]
This follows by checking the congruence modulo the coprime
polynomials \(S\) and \(X-1\).
Thus the circulant defined by the coefficients of \(H\) is the
inverse on the sum-zero subspace and vanishes on the constants.
If
\[
u=m^{-1},\qquad a=[X^0]H,\qquad b=[X^{m/2}]H,
\]
the entries of \(A_p^{-1}\) have exactly the form
\eqref{eq:Bsmall}. The entries involving zero follow from the
inverse of the displayed two-dimensional block. The equalities
at \(1,-1\) use \(5^{m/2}=-1\).

The polynomial Euclidean algorithm and resultant computation in
Appendix~\ref{app:code} give the exact data
\begin{equation}\label{eq:3data}
\Res(S,F)=2,\qquad u=1,\qquad a=1,\qquad b=0
\quad\text{in }\F_3.
\end{equation}
The inverse returned by the Euclidean algorithm is verified by a
separate multiplication and polynomial remainder calculation.
Equations \eqref{eq:resultant} and \eqref{eq:3data} give
\(\det A_p=2\), while
\[
\det C=\frac{1^2(1+0)}8=2\quad\text{in }\F_3.
\]
Finally, Proposition~\ref{prop:minor} yields
\[
D_{5469}=(-4)^{1823}\cdot2^3\cdot2=2
\quad\text{in }\F_3.
\]
Together with Section~\ref{sec:sun-prime}, this proves Theorem~\ref{thm:main}.


\section{Verification, provenance and historical scope}
\label{sec:scope}
Theorem~\ref{thm:cc} is an algebraic and real-variable proof for all
polynomial degrees; it does not depend on numerical verification.
Theorem~\ref{thm:main} is a complete exact counterexample to a universal
assertion, with its finite calculations reproduced below.
\subsection{Independent checks}
The structural certificate in Appendix~\ref{app:code} was executed
successfully. It was checked against independently implemented
calculations with different reductions:
\begin{center}
\begin{tabular}{@{}p{.54\textwidth}cc@{}}
\toprule
Calculation & modulo \(3\) & modulo \(1823\)\\
\midrule
Complementary minor and local polynomial certificate & \(2\) & \(1089\)\\
Direct original matrix, split under \(j\mapsto-j\) & \(1\cdot2\) & \(1054\cdot449\)\\
Literal \(5466\times5466\) matrix, direct elimination & \(2\) & ---\\
\bottomrule
\end{tabular}
\end{center}
In the second row, the two factors are determinants of blocks of
size \(2733\); their products are taken in the indicated fields.
The third row uses an integer Jacobi-symbol algorithm and direct
Gaussian elimination over \(\F_3\), without the Chinese remainder
or character reductions.
These checks support the certificate. The complete reproducible
proof route is the one supplied in the present manuscript.
No proof-assistant formalization or external referee assessment
is asserted.


\subsection{Relation to the supplied corpus}
The investigation used the supplied
\emph{Verifiable Mathematical Reasoning Corpus}, at commit
\texttt{c30d40119134a59eda6fa5cc6331f27c8e67d034}
\cite{Corpus}. Its work on root criteria, finite congruences, and
spectral reductions supplied directions for investigation.
Manuscript [077], the T96--T100 collection, led to Sun's determinant
problems; its T100 discussion concerns Conjecture~4.10(ii), a different
assertion. The finite congruence viewpoint of [089] was another useful
heuristic. No unverified corpus theorem is used as a premise here.
The targets and their original hypotheses were checked against the
primary papers \cite{ChasseCsordas2012,Sun2024}.

\subsection{What is proved and what is not certified}
The mathematical assertions are the original Chasse--Csordas
implication and the falsity of Sun's original universal assertion.
The stronger local-interval version is stated explicitly rather than
substituted for the original target.
For the first result, the forward measure formula was already proved
in Chasse's dissertation, and the decisive inverse-branch mechanism
is already present in Szeg\H{o}'s 1934 argument.
This manuscript supplies the full adaptation and cancellation-safe
recovery of the polynomial, but makes no claim to have invented that
general mechanism.

For the second result, the Chinese remainder theorem, character
diagonalization, polynomial Euclidean algorithms and the
complementary-minor identity are classical. Li--Yan already use
character diagonalization and Chinese remainder lifting in the same
determinant family. Here the additional calculations show that the
original truncated determinant is a unit at both prime factors of
$5469$, including the required nonzero complementary inverse minor.
This refutes the conjecture; it does not classify all integers in the
arithmetic progression, and no minimality claim for $5469$ is made.
Sun reports tests below $2000$ \cite[Remark~5.6]{Sun2024}; this
counterexample lies outside that range.

A bounded literature review checked the exact source conjectures,
Chasse's dissertation, the relevant classical rational-measure
argument, the dedicated Li--Yan paper, and searches for subsequent
resolutions and the specific counterexample. No identical later
resolution was located in that review. Such an absence is not a
certificate of priority or of the complete state of the literature.
The stronger requested success conditions therefore have the
following assessment.

\begin{center}\small
\begin{tabular}{@{}p{.06\textwidth}p{.84\textwidth}@{}}
\toprule
Gate & Assessment for the present manuscript\\
\midrule
G1 & Both targets have prior named and numbered source statements.\\
G2 & The finite-difference quotient, degree, measure law, original determinant, indices, and parameters are retained. Every change of representation is justified.\\
G3 & The source conjectures and inspected later treatments are documented; a universal absence of prior resolutions is not certified.\\
G4 & The polynomial implication is proved for all degrees and all hypotheses in its statement. The determinant assertion is refuted by an exact unconditional counterexample.\\
G5 & The arguments reach actual rigidity and nonvanishing. However, a historically unprecedented general principle is not established; Szeg\H{o}'s prior mechanism is material, and the determinant tools are classical.\\
G6 & No identical resolution was found in the bounded review. Search absence is not proof of firstness, and absolute priority is not asserted.\\
\bottomrule
\end{tabular}
\end{center}

The manuscript is a mathematical proof document for scrutiny. It does
not certify all historical or methodological success conditions, an
external referee's acceptance, or a higher significance ranking than
every previous manuscript in the corpus.
\appendix
\section{Complete exact certificate}\label{app:code}

The following program needs Python~3.8 or later and only its standard
library. It reads no files, downloads no data, and uses no floating-point
operations. Polynomial lists are in increasing order of degree.
The displayed code is embedded in this \texttt{.tex} file; no separate
source file is needed.

The routines \texttt{sub}, \texttt{mul} and \texttt{divrem} perform
polynomial arithmetic over the indicated prime field.
The extended Euclidean algorithm in \texttt{inverse\_mod} maintains
a coefficient of the input polynomial modulo the modulus.
The resultant routine uses, with \(R=A\bmod B\), \(d=\deg A\),
\(e=\deg B\), and \(r=\deg R\),
\[
\Res(A,B)=(-1)^{de}\operatorname{lc}(B)^{d-r}\Res(B,R),
\qquad \Res(A,b)=b^{\deg A}.
\]
This proves that the computed scalar is the resultant used in
\eqref{eq:resultant}. Inverse certificates are additionally checked
by multiplying them out. Primality of \(1823\) and \(911\) is
verified by exhaustive trial division through their integer square
roots, so no probabilistic primality test is involved.

The output is
\begin{verbatim}
n = 5469
mod p: Q = 1126 E = 1662 D = 1089
mod 3: det(Ap) = 2 a = 1 b = 0 inverse minor = 2 D = 2
Jacobi(D, 5469) = -1
\end{verbatim}
The program uses \(Q=\prod_{s=1}^{p-2}(-a_s)=-P\), so its
\(Q^3E\) is exactly \(-P^3E\) from \eqref{eq:prime-formula}.

\begin{lstlisting}
#!/usr/bin/env python3
"""Exact, standard-library-only certificate for Sun Conjecture 5.6(ii).

Every operation is integer arithmetic.  This verifies D={5,5}_5469 is
2 modulo 3 and 1089 modulo 1823, hence (D/5469)=-1, not 0.
The structural formulas are proved in Sections 5--7 of this manuscript.
"""

from math import isqrt


def trim(a):
    while a and not a[-1]:
        a.pop()
    return a


def sub(a, b, ell):
    c = [0] * max(len(a), len(b))
    for i, x in enumerate(a):
        c[i] += x
    for i, x in enumerate(b):
        c[i] -= x
    return trim([x % ell for x in c])


def mul(a, b, ell):
    if not a or not b:
        return []
    c = [0] * (len(a) + len(b) - 1)
    for i, x in enumerate(a):
        if x:
            for j, y in enumerate(b):
                c[i+j] += x*y
    return trim([x % ell for x in c])


def divrem(a, b, ell):
    assert b
    a = a[:]
    q = [0] * max(0, len(a)-len(b)+1)
    leading_inverse = pow(b[-1], -1, ell)
    while a and len(a) >= len(b):
        k = len(a)-len(b)
        t = a[-1] * leading_inverse % ell
        q[k] = t
        for j, x in enumerate(b):
            a[k+j] = (a[k+j]-t*x) % ell
        trim(a)
    return trim(q), a


def inverse_mod(a, modulus, ell):
    r0, r1 = modulus[:], divrem(a, modulus, ell)[1]
    t0, t1 = [], [1]
    while r1:
        q, r2 = divrem(r0, r1, ell)
        r0, r1 = r1, r2
        t0, t1 = t1, sub(t0, mul(q, t1, ell), ell)
    assert len(r0) == 1 and r0[0], "not invertible"
    t0 = [(x * pow(r0[0], -1, ell)) % ell for x in t0]
    return divrem(t0, modulus, ell)[1]


def resultant(a, b, ell):
    a, b = trim(a[:]), trim(b[:])
    res = 1
    while b:
        m, n = len(a)-1, len(b)-1
        if n == 0:
            return res * pow(b[0], m, ell) % ell
        r = divrem(a, b, ell)[1]
        if not r:
            return 0
        if m*n % 2:
            res = -res
        res = res * pow(b[-1], m-(len(r)-1), ell) % ell
        a, b = b, r
    return 0


def main():
    p = 1823
    n = 3*p
    assert n == 5469 and n % 60 == 9 and n > 69
    assert all(p % d for d in range(2, isqrt(p)+1))
    assert all(911 % d for d in range(2, isqrt(911)+1))
    m = p-1
    g = 5
    # m=2*911, so these two tests certify that g generates F_p^*.
    assert pow(g, m//2, p) != 1 and pow(g, m//911, p) != 1

    def chi(x):
        r = pow(x % p, m//2, p)
        return -1 if r == p-1 else r

    assert chi(5) == -1

    # Modulo p: coefficients of (X^2+5X+5)^((p-1)/2).
    a = [1]
    for _ in range(m//2):
        b = [0]*(len(a)+2)
        for i, x in enumerate(a):
            b[i] += 5*x
            b[i+1] += 5*x
            b[i+2] += x
        a = [x % p for x in b]
    assert len(a) == p and a[0] == p-1 and a[-1] == 1
    assert all(a[s] for s in range(1, p-1))
    Q = 1
    for s in range(1, p-1):
        Q = Q * (-a[s]) % p
    E = sum(pow(a[s], -1, p) for s in range(2, p-1, 2)) % p
    Dp = pow(Q, 3, p)*E % p
    assert (Q, E, Dp) == (1126, 1662, 1089)
    alpha = sum(pow(a[s], -1, p) for s in range(1, p-1)) % p
    beta = sum((-1)**s * pow(a[s], -1, p)
               for s in range(1, p-1)) % p
    assert (alpha, beta, (-Q) % p) == (921, 580, 697)
    assert pow(-4, p, p)*pow((-Q) % p, 3, p) % p == 641
    assert E*pow(4, -1, p) % p == 1327

    # Modulo 3: multiplicative circulant and its inverse off constants.
    # F_k=chi(g^(2k)+5g^k+5), S=1+X+...+X^(m-1).
    F = []
    z = 1
    for _ in range(m):
        F.append(chi(z*z+5*z+5) % 3)
        z = z*g % p
    assert z == 1 and sum(F) % 3 == 0
    S = [1]*m
    H = inverse_mod(F, S, 3)
    # Independently multiply the returned inverse, checking its certificate.
    assert divrem(mul(F, H, 3), S, 3)[1] == [1]
    H += [0]*(m-len(H))
    kshift = sum(H)*pow(m, -1, 3) % 3
    H = [(x-kshift) % 3 for x in H]
    assert sum(H) % 3 == 0
    cyclic_modulus = [2] + [0]*(m-1) + [1]
    target = [(1-pow(m, -1, 3)) % 3]
    target += [(-pow(m, -1, 3)) % 3]*(m-1)
    assert divrem(mul(F, H, 3), cyclic_modulus, 3)[1] == trim(target)
    aa, bb = H[0], H[m//2]
    det_Ap = m*resultant(S, F, 3) % 3
    u = pow(m, -1, 3)
    det_minor = u*u*(aa+bb)*pow(8, -1, 3) % 3
    D3 = pow(-4, p, 3)*pow(det_Ap, 3, 3)*det_minor % 3
    assert (det_Ap, aa, bb, det_minor, D3) == (2, 1, 0, 2, 2)

    # Legendre symbols at the two distinct prime factors of n.
    assert pow(Dp, m//2, p) == 1 and Dp == 33**2
    assert D3 == 2
    print("n =", n)
    print("mod p: Q =", Q, "E =", E, "D =", Dp)
    print("mod 3: det(Ap) =", det_Ap, "a =", aa, "b =", bb,
          "inverse minor =", det_minor, "D =", D3)
    print("Jacobi(D, 5469) = -1")


if __name__ == "__main__":
    main()
\end{lstlisting}

\begin{thebibliography}{9}
\raggedright

\bibitem{ChasseCsordas2012}
M. Chasse and G. Csordas,
\emph{Discrete analogues of the Laguerre inequalities and a conjecture
of I. Krasikov},
Annales des sciences math\'ematiques du Qu\'ebec
\textbf{36} (2012), no.~1, 77--91.
Definition~1.3, p.~78; Conjectures~2.21--2.22, p.~88.
\url{https://www.labmath.uqam.ca/~annales/volumes/36-1/PDF/077-091.pdf}.

\bibitem{ChasseThesis}
M. N. Chasse,
\emph{Linear preservers and entire functions with restricted zero loci},
Ph.D. dissertation, University of Hawai`i at Manoa, May 2011.
Theorem~88, pp.~45--46; Conjecture~89, p.~46.
\url{https://math.hawaii.edu/home/talks/thesis.pdf}.

\bibitem{Szego1934}
G. Szeg\H{o},
\emph{L\"osung der Aufgabe 107},
Jahresbericht der Deutschen Mathematiker-Vereinigung
\textbf{43} (1934), zweite Abteilung, 17--20.
Solution received 24 May 1931.
Authoritative digitization, scans 323--326:
\url{https://gdz.sub.uni-goettingen.de/id/PPN37721857X_0043}.

\bibitem{Sun2024}
Z.-W. Sun,
\emph{Problems and results on determinants involving Legendre symbols},
arXiv:2405.03626v8, 28 July 2024.
Definition (5.1), p.~22; Conjecture~5.6(ii) and Remark~5.6, p.~23.
\url{https://arxiv.org/abs/2405.03626v8}.

\bibitem{LiYan2026}
G. Li and X. Yan,
\emph{On Sun's conjectures for truncated Jacobi-symbol determinants
via supersingular elliptic curves},
arXiv:2608.08509v2, 23 October 2026.
Theorem~1.1(f), Remark~1.2, and Lemmas~2.1 and~2.4.
The author name in the inspected v2 PDF is Xiaojun Yan.
\url{https://arxiv.org/pdf/2608.08509v2}.

\bibitem{Corpus}
QEEDDD,
\emph{Verifiable Mathematical Reasoning Corpus},
research-draft repository, snapshot at commit
\texttt{c30d40119134a59eda6fa5cc6331f27c8e67d034},
6 November 2026.
Manuscripts [077] and [089] are contextual sources, not
logical dependencies.
\url{https://github.com/QEEDDD/Verifiable-Mathematical-Reasoning-Corpus}.

\end{thebibliography}
\end{document}
